\documentclass[11pt,letterpaper]{article}
\usepackage[margin=0.86in]{geometry}
\usepackage[T1]{fontenc}
\usepackage[utf8]{inputenc}
\usepackage{lmodern}
\usepackage{xcolor}
\usepackage{graphicx}
\usepackage{tabularx,booktabs,array,longtable}
\usepackage{enumitem}
\usepackage{hyperref}
\usepackage{fancyhdr}
\usepackage{lastpage}
\usepackage{titlesec}
\usepackage{amsmath,amssymb}
\usepackage{microtype}
\usepackage{needspace}
\usepackage{ragged2e}
\definecolor{navy}{HTML}{19334C}
\definecolor{teal}{HTML}{236B82}
\definecolor{soft}{HTML}{EFF5F7}
\definecolor{muted}{HTML}{52616B}
\hypersetup{colorlinks=true,linkcolor=teal,urlcolor=teal}
\titleformat{\section}{\Large\bfseries\color{navy}}{\thesection.}{0.6em}{}
\titleformat{\subsection}{\large\bfseries\color{navy}}{\thesubsection}{0.6em}{}
\titleformat{\subsubsection}{\normalsize\bfseries\color{navy}}{\thesubsubsection}{0.6em}{}
\titlespacing*{\section}{0pt}{15pt}{6pt}
\titlespacing*{\subsection}{0pt}{10pt}{4pt}
\setlength{\parindent}{0pt}
\setlength{\parskip}{4.5pt}
\setlength{\headheight}{15pt}
\renewcommand{\arraystretch}{1.16}
\setlist[itemize]{topsep=3pt,itemsep=2pt,leftmargin=1.4em}
\setlist[enumerate]{topsep=3pt,itemsep=2pt,leftmargin=1.5em}
\pagestyle{fancy}
\fancyhf{}
\fancyhead[L]{\small\color{muted} Applied Machine Learning \textbar{} Cornell Tech}
\fancyhead[R]{\small\color{muted} Human or AI?}
\fancyfoot[C]{\small\color{muted} Page \thepage\ of \pageref{LastPage}}
\renewcommand{\headrulewidth}{0.3pt}
\newcommand{\placeholder}[1]{\par\noindent\textcolor{teal}{\textit{[Your response: #1]}}\par\vspace{0.3em}}
\newcommand{\figplaceholder}[1]{\par\noindent\fcolorbox{teal}{soft}{\parbox[c][1.15in][c]{0.94\linewidth}{\centering\textcolor{muted}{Insert #1 here}}}\par}
\newcommand{\question}[2]{\needspace{4\baselineskip}\par\addvspace{0.45em}\noindent\textbf{#1}\quad #2\par\nobreak\vspace{0.1em}}
\newcommand{\smallnote}[1]{{\small\color{muted}#1}}
\newcolumntype{Y}{>{\RaggedRight\arraybackslash}X}
\begin{document}
\thispagestyle{empty}
\vspace*{0.2in}
{\color{teal}\bfseries CORNELL TECH \textbar{} APPLIED MACHINE LEARNING}\par
\vspace{0.15in}
{\Huge\bfseries\color{navy} Human or AI?}\par
\vspace{0.07in}
{\LARGE\color{navy} Kaggle Competition \&  Report Template}\par
\vspace{0.17in}

\textbf{Report Template: }\href{https://www.overleaf.com/project/6ae1fc85355dcf680137d898/share#f2777b20d3adeeecfac9096f78008e2f8daa0ff4be43de3b}{TEMPLATE}

\textbf{Github Repository Template: }\href{https://github.com/sekhari/human-or-ai-submissions}{REPO}

\textbf{Welcome to the competition!} Your task is to build a machine learning classifier that predicts whether a passage was written by a human or generated by an AI model. You will use a labeled training set, an additional unlabeled set, and a held-out Kaggle test set. The aim is not just a high leaderboard score: we want you to compare methods, design careful experiments, understand generalization, and explain what your results actually show.

You are encouraged to try several approaches. A well-tested logistic regression or support vector machine is just as legitimate as a neural network. The project is designed to be approachable on a laptop; a GPU is \emph{not required} for full credit.

\vspace{0.06in}
\begin{tabularx}{\linewidth}{@{}>{\bfseries}p{1.58in}Y@{}}
Team / members & Zhitong Liu (zl2358); Sorathorn Thongpitukthavorn (st2329); Keyan Miao (km2299) \\
Kaggle team name & team 35 \\
GitHub PR URL & [pending team approval] \\
\end{tabularx}
\vspace{0.06in}
\subsection*{Your data}
\begin{tabularx}{\linewidth}{@{}lrrY@{}}
\toprule
\textbf{File} & \textbf{Rows} & \textbf{Labels} & \textbf{Purpose} \\
\midrule
\texttt{train.csv} & 10,000 & Visible & Supervised model development: 5,000 human (0), 5,000 AI (1) \\
\texttt{unlabeled.csv} & 3,000 & Hidden & Optional input for semi-supervised methods; required experiment \\
\texttt{test.csv} & 3,000 & Hidden & Predict the probability of AI authorship for Kaggle \\
\texttt{train\_starter\_2000.csv} & 2,000 & Visible & Convenient subset of \texttt{train.csv} for quick experiments \\
\bottomrule
\end{tabularx}

\smallnote{The starter data is a subset, not 2,000 additional labeled examples. Passages are approximately 75--350 words, with mean length near 210 words.}

\subsection*{How the test set works}
The test set mixes familiar and unfamiliar writing domains and AI generators, and includes some adversarially modified AI-written passages. You will \emph{not} receive labels, source model names, domain labels, or difficulty memberships for individual test examples. Submit a probability in $[0,1]$ for each test ID, where 1 means AI-generated. Kaggle evaluates submissions using \textbf{area under the ROC curve (AUROC)}. The \textbf{public leaderboard} covers 900 examples (30\%); the \textbf{private leaderboard} covers 2,100 (70\%) and determines the final competition performance score.


\vfill
\newpage
\section{Competition rubric and submission instructions}
\begin{tabularx}{\linewidth}{@{}Y>{\raggedleft\arraybackslash}p{0.85in}@{}}
\toprule \textbf{Component} & \textbf{Points}\\ \midrule
Kaggle private leaderboard performance & 25 \\
Model comparison and selection & 15 \\
Effect of training data size & 10 \\
Learning from unlabeled data & 5 \\
Generalization and error analysis & 20 \\
Open-ended improvements and insights & 15 \\
Code submission and GitHub PR history & 10 \\
\midrule \textbf{Total} & \textbf{100} \\ \bottomrule
\end{tabularx}

\textbf{How to use this document.} Complete the response fields and tables below; insert plots where requested. Keep your answers concise, specific, and supported by actual experiments. Points appear in brackets beside each question; no separate writing-quality grade is assigned.

\textbf{How credit is awarded.} We grade the \emph{quality of the experimental method, evidence, and interpretation}, not whether a proposed idea succeeds. \textbf{Negative results can also earn full credit} when the comparison is valid, the results are reported accurately, and the conclusions are thoughtful. For each question, make clear what you tried, what you measured, and what you learned. Do not invent scores or results.

\textbf{Submit three items:} (1) final Kaggle submission; (2) this completed report as PDF; (3) one GitHub PR link in the course repository, with runnable code, instructions and meaningful commit history. Include the PR link on the title page. Confirm upload locations and deadline time on the course site. The template for the repository is \href{https://github.com/sekhari/human-or-ai-submissions}{here}

\textbf{Compute and models.} Use models and pipelines that you can run and reproduce. Laptop-friendly training is strongly recommended; use of more expensive methods is allowed but not required for full credit. Do not use external LLM APIs to directly label or predict the competition test passages. Use only the supplied labeled dataset for supervised labels; allowed pretrained weights and AI-assisted development should be disclosed.

\textbf{Validation.} For the required controlled experiments, reserve a fixed stratified \textbf{2,000-row validation set} from \texttt{train.csv}, leaving up to 8,000 rows for training. Use the \emph{same} held-out validation set for comparisons and do not tune on test labels. You may retrain your final Kaggle model on all 10,000 labeled examples after selecting your approach. Where possible, document seeds and avoid source/near-duplicate leakage.

\textbf{Kaggle grading (25 points).} Final performance will use the private AUROC, not leaderboard rank. The intended continuous rule is $25\,\mathrm{clip}((A-0.50)/(T-0.50),0,1)$, where $A$ is your private AUROC and $T$ is the announced full-credit threshold. A provisional value of $T=0.90$ is under consideration; the instructors will confirm the fixed threshold before launch. A reference word TF-IDF and logistic regression baseline achieves approximately 0.868 private AUROC.

\section{Team information and model summary}
\question{Q0.1 [0 points].}{List team members and their main contributions.}
Team members: Zhitong Liu (zl2358), Sorathorn Thongpitukthavorn (st2329), and Keyan Miao (km2299). Zhitong developed the character TF--IDF baseline, learning-curve and pseudo-labeling experiments, a wording-cue ablation, and the v3 holdout error analysis. Plum developed the multi-view models and v3 ensemble. Confirm the remaining individual contributions before submission.
\question{Q0.2 [0 points].}{Summarize your final approach in 3--5 sentences. What is the overall pipeline from raw text to probability?}
We turn each passage into several views: word and character patterns, a version that masks topic words, and measurements of writing style and formatting. Nine classifiers combine those views into an AI-authorship probability. A second stage adds high-confidence pseudo-labeled passages from the supplied unlabeled and test files, then blends its predictions with the first stage. The final output is a probability for each test ID; hidden test labels are never used for training.
\question{Q0.3 [0 points].}{Record your final scores (enter private score only once it is revealed).}
\begin{tabularx}{\linewidth}{@{}Yp{1.8in}@{}}\toprule Metric & Your score \\ \midrule Validation AUROC & 0.999551 (v3 holdout analogue) \\ Kaggle public AUROC & 0.98552 (full v3) \\ Kaggle private AUROC & [pending] \\ \bottomrule\end{tabularx}
\newpage
\section{Model comparison and selection (15 points)}
\textbf{Learning goal:} demonstrate sensible representation choices, comparisons of at least three approaches, appropriate hyperparameter selection, and evidence-based selection of the final classifier.

\question{Q1.1 [3 points].}{How did you preprocess and represent the text? Specify tokenization, normalization, word/character features or embeddings, and any feature extraction choices. Why?}
We made two copies of each passage: one with HTML break tags and extra whitespace removed for word and character patterns, and another retaining layout cues. TF--IDF gives more weight to patterns that are distinctive across passages. V3 uses groups of one to three words (short 3 to 5 character strings within and across words), and a topic-masked view that keeps grammar while hiding most content words. 
\newline

It also measures writing style, including vocabulary variety, repetition, sentence length, punctuation, and scores from a small trigram language model. These views can complement one another when the test topic changes. We kept layout cues separate because double newlines appear in 66\% of AI training passages but no human ones, an unreliable shortcut. Zhitong's independent baseline used only three- to five-character TF--IDF patterns; its vocabulary was learned from training rows alone.
\question{Q1.2 [5 points].}{Which \textbf{three or more} modeling approaches did you try? Keep the labeled validation partition fixed across models.}
\begin{tabularx}{\linewidth}{@{}p{1.24in}Yrrr@{}}\toprule
Model & Text representation & Train AUROC & Val. AUROC & Train time \\ \midrule
A: word logistic & One- and two-word patterns & 0.9932 & 0.9367 & 6.2 s \\
A2: Naive Bayes & One- and two-word patterns & 0.9790 & 0.8966 & 0.02 s \\
B: logistic & Word, character, and 72 style measures & 1.0000 & 0.9743 & 2.2 s \\
C: linear SVM & Same features as B & 1.0000 & 0.9766 & 8.2 s \\
D: LightGBM & Style measures and compressed word patterns & 1.0000 & 0.9822 & 0.6 s \\
E: blend of B/C/D & Combined word, character, and style views & 1.0000 & 0.9837 & 11.0 s \\
Zhitong: logistic & Character patterns only & 0.999897 & 0.968976 & 8.4 s \\
\bottomrule\end{tabularx}
All rows use the same stratified 8,000/2,000 split (seed 42). Model E combines B, C, and D with weights 20\%, 20\%, and 60\%. Zhitong's character-only model is an independent baseline. The later v3 pipeline reached 0.999551 AUROC in a holdout run, while its full-data Kaggle public score was 0.98552. Section 6 explains why the holdout result may be optimistic.
\question{Q1.3 [4 points].}{For at least \textbf{two} approaches, which hyperparameters and candidate values did you try? How did you select the final settings without using Kaggle test labels?}
\begin{tabularx}{\linewidth}{@{}p{1.18in}Yp{1.38in}@{}}\toprule Model & Parameters/values searched & Chosen setting \\ \midrule
B: logistic & $C\in\{0.5,1,2,3,5,8\}$; style scaling $\gamma\in\{0.10,0.15,0.25\}$ & $C=5$, $\gamma=0.15$ (inner CV 0.9694) \\
C: linear SVM & $C\in\{0.2,0.5,1.0,1.5,2.0\}$; sigmoid calibration & $C=1.0$ (inner CV 0.9727) \\
D: LightGBM & Leaves $\in\{15,31,63\}$; min. children $\in\{30,50\}$; rate $\in\{0.04,0.05\}$ & Leaves 31, min. 50, rate 0.05 (inner CV 0.9764) \\
Zhitong: logistic & Character $(3,5)$ fixed; $C\in\{0.1,1,10\}$ & $C=10$ \\
\bottomrule\end{tabularx}
For B, C, and D, Plum compared the settings above with five-fold cross-validation inside the 8,000 training rows (seed 42). The separate 2,000-row validation set was reserved for the model comparison; hidden test labels were never used. Zhitong used those same validation IDs to compare the character model's $C$ values. Its validation AUROCs were 0.933322, 0.961446, and 0.968976, so $C=10$ was selected. That score may be somewhat optimistic because the validation set helped choose $C$. The linear searches adjusted how strongly complexity was penalized; the tree search controlled model size and learning speed to limit overfitting.
\question{Q1.4 [3 points].}{Which approach did you choose for your final submission? Explain using validation performance, computational cost and/or stability. If you later changed the model, explain why.}
We first chose Model E because its 0.9837 validation AUROC beat the word-only baseline's 0.9367, and it trained in about 11 seconds on CPU. It combined word and character pattern models with a model based on writing style. It also kept 0.9804 AUROC when formatting was removed, whereas a layout-dependent tree model fell to 0.8109. An early Kaggle submission scored 0.9395 publicly, suggesting the fixed validation split missed some test differences. V2 added topic-masked patterns and a more stable blend, reaching 0.97446. We then submitted v3, which blends nine classifiers and adds examples given confident provisional labels by its first stage; it scored 0.98552 publicly. V3 is the team's current submission. Its 0.999551 local holdout AUROC came from a related 8,000/2,000 validation run, not the exact full-data submitted fit. Section 6 discusses this limitation.
\section{Effect of training data size (10 points)}
\textbf{Learning goal:} understand the effect of the number of labeled examples on fitting and generalization. This is an empirical exercise; no formal sample-complexity theory is required.

\question{Q2.1 [4 points].}{Choose one model configuration. Train it with 500, 1,000, 2,000, 5,000 and 8,000 examples sampled from your 8,000-row training partition. Keep the same validation set and hyperparameters.}
We held the character TF--IDF logistic model at $C=10$ and kept the same 2,000 validation passages. Each smaller training subset was sampled with class stratification and seed 42; we fitted a fresh vocabulary and classifier at each size.
\begin{tabularx}{\linewidth}{@{}rrrY@{}}\toprule Training rows & Training AUROC & Validation AUROC & Notes (optional) \\ \midrule
500 & 1.000000 & 0.904324 & Character TF--IDF, $C=10$ \\
1,000 & 1.000000 & 0.934335 & Same validation set \\
2,000 & 1.000000 & 0.947913 & Same validation set \\
5,000 & 0.999998 & 0.962446 & Same validation set \\
8,000 & 0.999897 & 0.968976 & Same validation set \\
\bottomrule\end{tabularx}
\question{Q2.2 [3 points].}{Insert a plot of training and validation AUROC versus training-set size.}
\begin{center}\includegraphics[width=0.85\linewidth]{figures/baseline_learning_curve.pdf}\end{center}
\question{Q2.3 [3 points].}{What does the curve suggest about overfitting, underfitting, and the value of more labeled examples?}
Training AUROC stays near 1.0, while validation AUROC rises by 0.0647 between 500 and 8,000 passages. The gap is largest at 500, consistent with overfitting when the model has fewer labeled examples. The training scores show little evidence of underfitting. The gains shrink as the training set grows: the final 3,000 passages add about 0.0065 AUROC. More labels may still help, but this single split and one draw per size cannot tell us how much.
\section{Learning from unlabeled data (5 points)}
\textbf{Learning goal:} test whether the additional 3,000 unlabeled passages help.
\question{Q3.1 [2 points].}{Describe one method that uses \texttt{unlabeled.csv} (e.g., self-training, confidence-based pseudo-labeling, clustering, or feature learning). How did you prevent validation-label leakage?}
We trained Zhitong's character baseline on the 8,000 labeled training rows and scored the unlabeled passages. We dropped duplicate unlabeled texts and any exact matches to labeled texts, without consulting validation labels. We assigned pseudo-labels of human for $p\leq0.05$ and AI for $p\geq0.95$, retaining 540 and 572 passages respectively. We refit the same pipeline on all 8,000 training rows plus the 1,112 pseudo-labeled rows. Validation labels were used only to score this comparison. Plum's final v3 uses a separate pseudo-labeling configuration, so these two results should not be conflated.
\question{Q3.2 [2 points].}{Compare against a supervised baseline under the same held-out validation protocol.}
\begin{tabularx}{\linewidth}{@{}Yp{1.5in}@{}}\toprule Method & Validation AUROC \\ \midrule Supervised only (character baseline) & 0.968976 \\ With 1,112 pseudo-labeled rows & 0.968447 \\ \bottomrule\end{tabularx}
\question{Q3.3 [1 point].}{Did it help? Give a concrete, evidence-based explanation (including possible noise in pseudo-labels).}
The character-baseline pseudo-labeling experiment did not improve validation AUROC; it fell by 0.000529. Confidence does not guarantee a correct label, and adding the teacher's predictions can reinforce its mistakes. We do not know the true labels of the added passages, so label noise is a plausible explanation. One split and one threshold pair do not establish a general effect, and they do not determine whether the different self-training stage in v3 helps.
\newpage
\section{Generalization and error analysis (20 points)}
\textbf{Learning goal:} characterize how well the selected model generalizes and where it fails. Use \textbf{labeled validation data}; hidden test labels and hidden domain/generator information are unavailable.

For this section, we evaluated Plum's v3 architecture on a stratified seed-42 split of the 10,000 labeled passages: 8,000 for fitting and 2,000 held out. The v3 language-model features and human-reference genre statistics were fitted using only the 8,000 training labels. Its second stage pseudo-labeled passages from \texttt{unlabeled.csv} and \texttt{test.csv} without using validation labels. This CPU validation run fitted each specialist once; the Kaggle submission instead bags five fits trained from all 10,000 labeled rows. The numbers below therefore evaluate the v3 pipeline on a clean holdout, not the exact submitted fit.

\question{Q4.1 [5 points].}{Compare training and validation AUROC for your final selected model. Is there a gap? What explanations are plausible?}
The full two-stage v3 validation run reached 1.000000 in-sample training AUROC and 0.999551 held-out AUROC, a gap of 0.000449. This small within-dataset gap does not establish robust generalization: the formatting-free stage-one branch scored 0.987841 on the same holdout. In the labeled training data, 66.0\% of AI passages but 0\% of human passages contain a double newline, so paragraph structure is an unusually strong shortcut. The full-data v3 submission scored 0.98552 on Kaggle's public test subset, further suggesting that the random holdout is optimistic for the shifted test distribution. The public score is not a passage-level error analysis, and the private score remains unknown.
\question{Q4.2 [5 points].}{Evaluate performance over at least three text-length bins (define the bins), or another observable subgroup available in your data. Explain why you chose the grouping.}
We used Plum's pre-existing whitespace-delimited word-count bins to check whether passage length changes performance. All three contain several hundred validation passages and both classes.
\begin{tabularx}{\linewidth}{@{}Yrr@{}}\toprule Subgroup / bin & Count & Validation AUROC \\ \midrule
75--165 words & 578 & 0.999661 \\
166--245 words & 832 & 0.999147 \\
246--350 words & 590 & 0.999942 \\
\bottomrule\end{tabularx}
The groups contain 260/318, 481/351, and 259/331 human/AI passages respectively. At a diagnostic threshold of 0.49, their false-positive/false-negative counts are 4/2, 2/8, and 0/2. The middle bin has eight of the twelve false negatives; with so few mistakes, this is a diagnostic observation rather than strong evidence that medium-length text is intrinsically harder.
\question{Q4.3 [5 points].}{Analyze at least \textbf{five false positives and five false negatives} from your held-out validation set. Provide example IDs, predicted probabilities, and short failure hypotheses. Do not paste full passages unless essential.}
At the conventional 0.50 threshold, v3 has only four false positives and twelve false negatives on this holdout. We use 0.49 for the diagnostic table so that at least five errors of each type can be examined; it yields six FP and twelve FN. These are hypotheses from brief text inspection, not verified causes.
\begin{tabularx}{\linewidth}{@{}p{0.64in}p{0.70in}p{0.98in}Y@{}}\toprule ID & Error type & P(AI) & Possible reason \\ \midrule
7767 & FP & 0.6644 & Compact biography and credentials resemble a polished generated summary. \\
4062 & FP & 0.5646 & Song-release and chart facts form an orderly encyclopedia-style summary. \\
3660 & FP & 0.5634 & A disambiguation list has repeated names and line breaks. \\
5079 & FP & 0.5619 & League-season overview repeats formal descriptive phrasing. \\
3571 & FP & 0.4986 & Short film review repeats evaluative language and the word ``movie.'' \\
1362 & FN & 0.0535 & AI biography uses dates, proper nouns, and historical detail. \\
9298 & FN & 0.1482 & AI celebrity biography includes many names and dated events. \\
4519 & FN & 0.1503 & Lowercase first-person forum style resembles informal human writing. \\
1834 & FN & 0.2234 & AI country history uses an encyclopedia-like sequence of dates. \\
5498 & FN & 0.2277 & Ingredient quantities and line breaks resemble a human recipe. \\
\bottomrule\end{tabularx}
\smallnote{FP means human predicted AI; FN means AI predicted human. The threshold is 0.49.}
\question{Q4.4 [5 points].}{What recurring failure patterns did you identify, and how might you address one of them? Distinguish observed patterns from speculation.}
The selected errors show a genre-confusion pattern: several human encyclopedia-style passages are scored as AI, while several AI biographies and one AI recipe are scored as human. One AI forum-style passage is also a false negative, so no single genre explains every mistake. These ten examples do not measure the prevalence of each failure mode. A more consequential observed vulnerability is formatting dependence: within stage one, the full blend scored 0.999561 held-out AUROC versus 0.987841 for its formatting-free branch, while the public test score for the final v3 submission is lower than the random-holdout result. We would test a smaller structural-branch weight using held-out genres and line-break perturbations, then compare it with the current blend before changing the submission. That experiment could show whether the apparent formatting advantage survives distribution shift; the present data do not establish that it will improve private AUROC.
\newpage
\section{Open-ended improvements and insights (15 points)}
\textbf{Learning goal:} propose and test one improvement beyond the required baseline comparisons. This can be simple: a targeted feature change, alternate regularization scheme, ensemble, different representation, augmentation, or training approach. Sophistication alone does not earn credit.

\question{Q5.1 [3 points].}{What specific question or hypothesis motivated your improvement? What weakness or opportunity did you observe?}
The \href{https://en.wikipedia.org/wiki/Wikipedia:Signs_of_AI_writing}{Wikipedia field guide on signs of AI writing} describes rhetorical patterns that sometimes occur in AI-generated text, including negative parallelism and lists of three. I asked whether explicit counts of such patterns would add useful information beyond the character TF--IDF features in my baseline. The hypothesis was that a small set of interpretable structural cues could improve validation AUROC when combined with the existing classifier. These patterns are common in human writing too, and the guide's examples may reflect Wikipedia-specific editing practices, so I treated them as testable cues rather than rules for identifying authorship.
\question{Q5.2 [4 points].}{What exactly did you change? Describe the baseline, the new approach and what you held fixed for a fair comparison.}
The baseline was the selected character 3--5-gram TF--IDF logistic model with $C=10$. For the augmented version, I counted matches of \texttt{not ... but}, \texttt{rather than}, and an approximate three-item-list regex in each passage, divided each count by passage length per 100 words, standardized the three features using training rows only, multiplied them by 0.1, and appended them to the TF--IDF matrix. Both models used the same 8,000 training and 2,000 validation rows from seed 42, the same $C=10$, and the same AUROC metric. The cue patterns were applied to raw text; no validation labels entered feature fitting. The list regex is an approximation, not a grammatical parser.
\question{Q5.3 [4 points].}{What were the measured results?}
\begin{tabularx}{\linewidth}{@{}YrrY@{}}\toprule Configuration & Val. AUROC & Runtime & Relevant difference \\ \midrule Char TF--IDF baseline & 0.968976 & 8.07 s & 180 errors at 0.5 \\ + wording cues & 0.968779 & 10.42 s & 181 errors at 0.5 \\ \bottomrule\end{tabularx}
\question{Q5.4 [4 points].}{What did you learn? Did the results support the hypothesis? What limits the conclusion, and what would you try next?}
The explicit cues did not support the improvement hypothesis: validation AUROC fell from 0.968976 to 0.968779 (a difference of $-0.000197$), while 0.5-threshold errors rose from 180 to 181. In the training partition, \texttt{not ... but} appeared in 6.1\% of human and 7.8\% of AI passages, and \texttt{rather than} in 1.5\% and 2.9\%; the approximate three-item-list pattern appeared in 53.0\% of human and 45.9\% of AI passages. These overlapping rates help explain why the cues do not reliably identify authorship. The test covers only three approximate regex features, one fixed weight, one $C$, and one validation split. It does not rule out all stylistic features. I would next examine errors by passage type and test better-defined linguistic cues within training-only cross-validation before evaluating a revised model on the held-out set. We kept the character baseline simple; the team's final v3 submission was chosen from separate experiments.
\section{Code submission, reproducibility and AI assistance (10 points)}
\textbf{Submission format:} Submit a \textbf{GitHub PR URL} against the course repository. We will inspect the final code and its development history, not just the final snapshot. Do not squash away meaningful intermediate commits before review. Template \href{https://github.com/sekhari/human-or-ai-submissions}{here}.

\question{Q6.1 [5 points].}{Provide the PR URL, the entrypoint for training/evaluation, and the command used to generate \texttt{submission.csv}. State dependencies and random seeds.}
\placeholder{PR URL; commands; environment; output path}
\question{Q6.2 [3 points].}{Give two or three examples of changes you made during development. Link to relevant commits or PR discussion and briefly explain what the experiments taught you.}
\placeholder{Commit links and brief notes}
\question{Q6.3 [2 points].}{How did you use AI tools (such as Claude, ChatGPT or Copilot) for code, debugging, analysis or writing? What did you verify or modify yourself?}
\placeholder{Short, honest disclosure. Full prompt transcripts are not required.}

\textbf{AI policy.} You may use AI assistants for coding, debugging, explanations, brainstorming, and editing. We expect each team to exercise its own judgment, run experiments, check outputs, and understand its submitted approach. Sending the entire assignment to an AI tool and submitting its unverified output without meaningful involvement does not satisfy the objectives of this project. We will use the PR history, reproducible experiments and report together as evidence of engagement. \emph{Commit counts and the presence of AI-generated code are not themselves grading criteria.} Do not fabricate results or retrieve hidden test labels.


\section*{Final submission checklist}
\begin{itemize}[label=$\square$]
\item Final probabilities submitted to Kaggle for every \texttt{test.csv} ID.
\item Completed report, with all required answers, tables, figures and sources, submitted as PDF.
\item GitHub PR linked on the report title page and in the course submission portal.
\item Code and dependencies recreate final predictions (within reasonable numerical tolerance).
\item Relevant experiments and seed/configuration details are recorded.
\item AI-use disclosure is complete; hidden labels were not retrieved externally.
\end{itemize}


\end{document}
